\documentclass[11pt]{article}
\usepackage[margin=2.5cm]{geometry}
\usepackage{amsmath,amssymb}
\usepackage{booktabs}
\usepackage{enumitem}
\usepackage[hidelinks]{hyperref}

\title{Lessons from the Real-Data Track\\
\large fracres meets ds002778: implementation, failures, and validations}
\author{Working note --- fracres project (D.~R.~Chin)}
\date{15 September 2026}

\begin{document}
\maketitle

\begin{abstract}
\noindent This note records what was learned while taking fracres from a
phantom generator to an instrument that loads real EEG, reconstructs
held-out channels, and infers generative mechanisms (ROADMAP ``Next 1'',
commits \texttt{92b473a}, \texttt{d96da46}, \texttt{8abc81e}). It covers the
three scientific findings, the failures encountered en route and how each was
resolved, the validation patterns that proved trustworthy, and the open
threads. The failures are recorded deliberately: several were cheap to make
and would be cheap to make again.
\end{abstract}

\section{What was built}

Three modules over the UC San Diego resting-state Parkinson's dataset
(OpenNeuro \texttt{ds002778}: 15 PD patients $\times$ medication OFF/ON
sessions, 16 controls; BioSemi 32-channel, 512\,Hz, ${\sim}200$\,s runs;
571\,MB mirrored locally from the anonymous S3 bucket):

\begin{enumerate}[leftmargin=1.6em]
  \item \textbf{Ingestion} (\texttt{fracres.data}): BIDS- and session-aware
        loading (\texttt{ses-off}/\texttt{ses-on}/\texttt{ses-hc}),
        participants metadata, EXG/status channel dropping, standard-montage
        geometry, optional resampling, JAX hand-off. \texttt{mne} is a lazy
        optional extra.
  \item \textbf{Forward fit} (\texttt{fracres.forward\_fit}): average
        re-reference + z-score, drive the frozen reservoir with a real channel,
        closed-form ridge reconstruction of held-out channels (temporal split
        + washout), per-channel held-out correlation/nMSE and DFA-$H$ /
        spectral-$\beta$ comparison.
  \item \textbf{Inverse problem} (\texttt{fracres.inverse}): weighted Euclidean
        objective on the (DFA-$H$, spectral-$\beta$) metric vector;
        declarative $(\alpha, H, \lambda)$ grid search with common random
        numbers, seed-ensembled metrics, length-matched simulation, and an
        \texttt{fbm} drive regime.
\end{enumerate}

\section{Scientific findings}

\paragraph{1. The OFF/ON asymmetry (forward fit).}
Driven by Fz of \texttt{sub-pd3}, the same $(\alpha, H)$ phantom reconstructs
seven channels at mean correlation $0.52$ OFF medication but only $0.19$ ON;
the data's statistics shift between sessions (spectral $\beta$ up to
$1.5$--$1.8$, DFA-$H$ to $1.3$--$1.45$). No single mechanism matches both
medication states --- the inverse problem is empirically necessary, not merely
planned.

\paragraph{2. The fGn regime mismatch, and its resolution.}
Raw scalp EEG is fBm-\emph{like}: $\beta \approx 1.8$, DFA-$H \approx 1.3$
(nonstationary). A \emph{contractive} reservoir driven by stationary fGn
cannot produce $\beta > 1$ --- the first grid sweep bottomed out at objective
$\approx 1.58$ with every candidate near DFA-$H \approx 0.7$, $\beta \approx
0.4$. The failure was itself the diagnostic: the drive must be integrated.
With \texttt{drive\_kind="fbm"} (z-scored cumulative fGn), the stable
reservoir passes the drive's low-frequency power law through, so
\emph{anti-persistent} $H$ reaches the EEG regime exactly:
$\beta = 2H+1$, DFA-$H = H+1$.

\paragraph{3. The medication contrast (inverse problem).}
Best-fit mechanism for \texttt{sub-pd3} Cz (60-point grid, 2 seeds):
OFF $\to (\alpha{=}0.9, H{=}0.7, \lambda{=}1)$, objective $0.068$;
ON $\to (\alpha{=}0.9, H{=}0.6, \lambda{=}1)$, objective $0.008$ (DFA-$H$
$1.398$ vs.\ data $1.403$; $\beta$ $1.94$ vs.\ $1.94$). Dopamine leaves the
fractional \emph{order} unchanged and \emph{whitens the effective drive}
($H\colon 0.7 \to 0.6$). Caveats: one subject, $\alpha$ at the grid boundary,
$H$ spacing $0.1$. A hypothesis, not a result.

\section{Failures encountered and how they were resolved}

Recorded most-recent-first within each category; each entry names the failure,
the signal that exposed it, and the fix.

\subsection{Latent code defects}

\begin{description}[leftmargin=1.6em,style=nextline]
  \item[\texttt{signal\_metrics} argument misrouting.]
  The bundle passed \texttt{n\_segments=8} into \texttt{spectral\_exponent}'s
  \texttt{f\_max} slot, so every bundled $\beta$ was silently a \emph{full-band}
  fit ($f_{\max}=8$ cycles/sample --- above Nyquist) --- precisely the bias the
  $f_{\max}$ restriction exists to prevent. The existing test tolerated it
  because its $0.15$ agreement slack absorbed the underestimate. \emph{Fix:}
  signature changed to \texttt{signal\_metrics(x, f\_max=0.1)}; regression test
  pins exact equality with \texttt{spectral\_exponent(x)} and the direction of
  the band-restriction effect. Estimator recovery improved visibly
  ($H{=}0.8$: $\beta$ $0.555$ vs.\ true $0.6$). \emph{Lesson:} positional
  arguments across similarly-shaped signatures are a standing hazard; bundle
  functions should forward by keyword, and loose agreement tests hide
  misrouted arguments.

  \item[Suite-order-dependent test failure (global \texttt{x64}).]
  A new regression test passed standalone but failed in the full suite:
  other test modules enable \texttt{jax\_enable\_x64} globally at import,
  and the Davies--Harte fGn differs under float32/float64 --- float32 adds a
  high-frequency noise floor that flattens the spectrum enough to flip an
  inequality ($\beta_{0.1} = 0.587$ vs.\ full-band $0.401$ in float32;
  $0.620$ vs.\ $0.622$ in float64). \emph{Fix:} the directional assertion now
  uses a prescribed-spectrum signal (exactly $f^{-0.6}$ below $f{=}0.1$, flat
  above, random phases), which is strictly ordered in any precision.
  \emph{Lesson:} never write order-dependent assertions on stochastically
  generated signals when global precision state varies; construct the
  property you want to test directly.

  \item[Asserting an invariant in the wrong domain.]
  My own preprocessing test asserted that average-referenced channels sum to
  zero \emph{after} per-channel z-scoring --- false, because z-scoring
  rescales channels unequally. \emph{Fix:} assert the zero-sum property with
  \texttt{zscore=False} and the unit-variance property separately.
  \emph{Lesson:} state which processing stage each invariant belongs to;
  composing transforms composes their side conditions.
\end{description}

\subsection{Wrong assumptions about the dataset}

All three of these would have been avoided by inspecting one subject's files
before writing the loader --- they were caught by tests against the real
mirror instead:

\begin{itemize}[leftmargin=1.6em]
  \item \textbf{Session level.} ds002778 is not
        \texttt{sub-*/eeg/} but \texttt{sub-*/ses-*/eeg/}, and PD subjects
        have \emph{two} sessions (\texttt{ses-off}/\texttt{ses-on},
        medication state) --- which turned out to be the most scientifically
        valuable axis of the whole track. \texttt{recording\_path} now
        resolves sessions, auto-selecting a single session and requiring an
        explicit choice when there are several.
  \item \textbf{Channel count.} Assumed BioSemi 64 (from the EXG1--8
        discussion in the dataset comments); the cap is 32 scalp channels.
        Default montage corrected to \texttt{biosemi32} (and the mne naming
        has no underscore --- \texttt{biosemi\_64} fails with a long
        allowed-values error).
  \item \textbf{Participants schema.} No \texttt{Group} column; group comes
        from the subject-label prefix. The parser now falls back to label
        inference.
\end{itemize}

\subsection{Test-calibration failures}

\begin{itemize}[leftmargin=1.6em]
  \item \textbf{Arbitrary margins.} The inverse-problem recovery test first
        demanded every non-true grid point score $> 0.05$; the runner-up
        scored $0.039$ --- legitimately, since $\alpha$ is the soft direction
        at finite $T$. Replaced with \emph{structurally motivated}
        assertions: the true point scores $\approx 0$ (exact replay via
        common random numbers), wrong-$H$ points lose by $\geq 0.3$
        (the sharp direction). Test now encodes an identifiability claim, not
        a magic number.
  \item \textbf{Mental arithmetic in expected values.} A weighted-distance
        test expected $0.4$ where the definition gives $\sqrt{4 \cdot 0.1^2}
        = 0.2$; the module was right. Weights multiply squared deviations ---
        document the convention at the signature.
\end{itemize}

\subsection{Environment and access}

\begin{itemize}[leftmargin=1.6em]
  \item OpenNeuro needs neither DataLad nor the AWS CLI: the
        \texttt{openneuro.org} S3 bucket allows anonymous \texttt{ListBucket}
        and \texttt{GET}, so a 60-line stdlib script mirrors a dataset
        (resumable, size-verified). \texttt{mne.export} to EDF additionally
        needs \texttt{edfio}; BDF export is unsupported, so synthetic-test
        fixtures are EDF.
\end{itemize}

\section{Validation patterns that earned trust}

\begin{enumerate}[leftmargin=1.6em]
  \item \textbf{Exact replay (common random numbers).} Sharing model/drive
        seeds across grid points makes the true mechanism score exactly $0$
        in a recovery test --- a deterministic, epsilon-free identifiability
        check that doubles as variance reduction in the search itself.
  \item \textbf{Expressible vs.\ unpredictable channels.} The forward-fit
        unit test pairs a lagged/smoothed target (must reconstruct) with an
        independent-noise channel (must not), so the pipeline is tested at
        both ends of its operating range.
  \item \textbf{Opt-in integration tests.} Real-data tests skip unless the
        local mirror exists (same pattern as the hpfracc cross-reference):
        CI stays data-free, but every structural assumption is verified
        against the real recordings on a developer machine.
  \item \textbf{Metrics with physical priors.} Knowing EEG \emph{should} sit
        at $\beta \in [1,2]$ turned the inverse problem's large residual from
        a disappointing number into a mechanism diagnosis (Section~2,
        finding 2). Sanity ranges are part of the instrument.
\end{enumerate}

\section{Open threads}

\begin{itemize}[leftmargin=1.6em]
  \item Grid refinement: $\alpha \to 0.95/1.0$ (best-fit is at the boundary),
        $H$ at $0.05$ spacing, more seeds for MLE-tight error bars.
  \item All-subject $\times$ all-channel OFF/ON sweep --- the population-level
        version of finding 3, and the seed of the phantom-vs-real-EEG study
        (ROADMAP Next 5).
  \item Per-node $(\alpha, \lambda)$ (Next 2) as the mechanism question moves
        from ``which order'' to ``which distribution of orders''.
  \item The \texttt{fbm} drive regime currently lives in
        \texttt{fracres.inverse}; promoting it to \texttt{fracres.drivers}
        would let every example and the config system use it.
  \item Data-use obligations for any publication: cite Jackson et al.\ 2019
        (\emph{eNeuro}), Swann et al.\ 2015 (\emph{Ann.\ Neurol.}), George et
        al.\ 2013 (\emph{NeuroImage Clin.}); email the curators
        pre-submission; no PD/HC classifier on this $n$ (curators' explicit
        request --- our use is forward modelling and LRD statistics).
\end{itemize}

\end{document}
